% --- Archon physics formalization source begin ---
% archon:chemistry
% archon:covers IChO2026Problems/problem_icho_2026_t1_a6.lean
% archon:problem-id icho_2026_t1
% archon:part-id T1-A6
% archon:previous-part icho_2026_t1_a4
% archon:previous-part-policy derive_in_answer_blind_run_or_use_problem_stated_fallback
% archon:previous-part icho_2026_t1_a5
% archon:previous-part-policy derive_in_answer_blind_run_or_use_problem_stated_fallback

\chapter{Icho Chemistry problem icho\_2026\_t1\_a6}
\label{ch:IChO2026Problems_problem_icho_2026_t1_a6}

\paragraph{Problem source.}
\#\# Shared official problem context
T1. A journey through time: The secrets of Avicenna 7\% A chemist found herself in the 11 th Century, where she met the renowned physician Avicenna. In Avicenna's laboratory, she discovered a blue elixir and a yellow mysterious stone. She brought them back to 2026. Determined to uncover their composition, she began her research. Part 1. The blue elixir At first, the chemist compiled a table listing the names of the only four plants which she saw in Avicenna’s lab that could have been used to prepare the elixir, and the compounds that can be extracted from these plants. Plants Compounds that can be extracted Zingiber 1 (C11H14O3) 2 (C10H18O) 3 (C10H18O) Hypericum 4 (C6H12O) 5 (C10H12O2) 6 (C10H18O) Chamomilla 7 (C14H16) 8 (C15H24) 3 (C10H18O) Artemisia 9 (C10H16O) 10 (C10H18O) 3 (C10H18O) Using chromatography, she determined that the elixir consists of four different substances (X, Y, Z and W). The chemist then discovered that X can isomerise into Y in an acidic medium and that Y has a plane of symmetry. 1.1 Identify X and Y. Chromatography also revealed it was substance Z that gave the elixir its blue colour. Z is a derivative of the well-known substance A, which has two perpendicular planes of symmetry. When heated, A isomerises into aromatic compound B, which has three mutually perpendicular planes of symmetry. When substance Z is chlorinated, a pair of peaks with an intensity ratio of 3:1 is observed in the mass spectrum at m/z 218 and 220, respectively. 1.2 Identify Z and draw the structures of substances A and B. The chemist extracted W from the elixir. When W was added to an aqueous Fe3+ solution, a characteristic colour change was observed. Mass spectrometric analysis of W revealed the intensity ratio of [M]+ ∶[M + 1]+= 9:1, where M is the molecular ion. Assume carbon consists exclusively of 12C and 13C, the natural isotopic abundance of 12C is 98.9\%, and all other elements are monoisotopic. 1.3 Determine the number of carbon atoms (𝑛) from the mass spec data and identify W. Show your calculations. Part 2. The mysterious stone Avicenna had told her the stone was a stoichiometric compound. She dissolved it in dilute nitric acid. After that, she adjusted the pH of the solution to \textasciitilde{}4 and added NaF. This led to the formation of precipitate C ⋅𝑥H2O, which contained 39.16\% water by mass. Anhydrous C reacts with an excess of NaF to yield compound D, which contains 32.85\% sodium and 12.85\% of metal Q by mass, and is used in the industrial production of Q. 1.4 Identify metal Q and compounds C ⋅𝑥H2O and D. Avicenna told her that he found the stone in a coal deposit. The stone contains the anion of a highly symmetrical acid F. Reaction of F with P2O5 gives binary compound G, which contains 49.98\% of oxygen by mass and has a three-fold symmetry axis. F can be prepared by oxidising compound E with potassium permanganate in an acidic solution. E contains 11.18\% of hydrogen by mass and has a six-fold symmetry axis. 1.5 Draw the structures of E, F, and G. Then, she analysed the stone thermogravimetrically in open air. A 10.00 g sample began to lose mass at approximately 100 ∘C and stopped at 5.75 g at 200 ∘C. A second drop in mass was observed at 400 ∘C, with a final mass of 1.50 g of compound H remaining constant at higher temperatures.

\#\# Current subquestion T1-A6
Determine the chemical formulae of the stone and compound H using thermogravimetric data.

\paragraph{Shared problem context.}
T1. A journey through time: The secrets of Avicenna 7\% A chemist found herself in the 11 th Century, where she met the renowned physician Avicenna. In Avicenna's laboratory, she discovered a blue elixir and a yellow mysterious stone. She brought them back to 2026. Determined to uncover their composition, she began her research. Part 1. The blue elixir At first, the chemist compiled a table listing the names of the only four plants which she saw in Avicenna’s lab that could have been used to prepare the elixir, and the compounds that can be extracted from these plants. Plants Compounds that can be extracted Zingiber 1 (C11H14O3) 2 (C10H18O) 3 (C10H18O) Hypericum 4 (C6H12O) 5 (C10H12O2) 6 (C10H18O) Chamomilla 7 (C14H16) 8 (C15H24) 3 (C10H18O) Artemisia 9 (C10H16O) 10 (C10H18O) 3 (C10H18O) Using chromatography, she determined that the elixir consists of four different substances (X, Y, Z and W). The chemist then discovered that X can isomerise into Y in an acidic medium and that Y has a plane of symmetry. 1.1 Identify X and Y. Chromatography also revealed it was substance Z that gave the elixir its blue colour. Z is a derivative of the well-known substance A, which has two perpendicular planes of symmetry. When heated, A isomerises into aromatic compound B, which has three mutually perpendicular planes of symmetry. When substance Z is chlorinated, a pair of peaks with an intensity ratio of 3:1 is observed in the mass spectrum at m/z 218 and 220, respectively. 1.2 Identify Z and draw the structures of substances A and B. The chemist extracted W from the elixir. When W was added to an aqueous Fe3+ solution, a characteristic colour change was observed. Mass spectrometric analysis of W revealed the intensity ratio of [M]+ ∶[M + 1]+= 9:1, where M is the molecular ion. Assume carbon consists exclusively of 12C and 13C, the natural isotopic abundance of 12C is 98.9\%, and all other elements are monoisotopic. 1.3 Determine the number of carbon atoms (𝑛) from the mass spec data and identify W. Show your calculations. Part 2. The mysterious stone Avicenna had told her the stone was a stoichiometric compound. She dissolved it in dilute nitric acid. After that, she adjusted the pH of the solution to \textasciitilde{}4 and added NaF. This led to the formation of precipitate C ⋅𝑥H2O, which contained 39.16\% water by mass. Anhydrous C reacts with an excess of NaF to yield compound D, which contains 32.85\% sodium and 12.85\% of metal Q by mass, and is used in the industrial production of Q. 1.4 Identify metal Q and compounds C ⋅𝑥H2O and D. Avicenna told her that he found the stone in a coal deposit. The stone contains the anion of a highly symmetrical acid F. Reaction of F with P2O5 gives binary compound G, which contains 49.98\% of oxygen by mass and has a three-fold symmetry axis. F can be prepared by oxidising compound E with potassium permanganate in an acidic solution. E contains 11.18\% of hydrogen by mass and has a six-fold symmetry axis. 1.5 Draw the structures of E, F, and G. Then, she analysed the stone thermogravimetrically in open air. A 10.00 g sample began to lose mass at approximately 100 ∘C and stopped at 5.75 g at 200 ∘C. A second drop in mass was observed at 400 ∘C, with a final mass of 1.50 g of compound H remaining constant at higher temperatures.

\paragraph{Current subquestion.}
Determine the chemical formulae of the stone and compound H using thermogravimetric data.

\paragraph{Figure/image paths.}
\begin{itemize}
\item icho\_2026\_source/image/T1\_page-4.png
\item icho\_2026\_source/image/T1\_page-3.png
\end{itemize}

\paragraph{Reusable previous-part conclusions.}
\begin{itemize}
\item Source T1-A4. Question: Identify metal Q and compounds C ⋅𝑥H2O and D. Policy: derive\_in\_answer\_blind\_run\_or\_use\_problem\_stated\_fallback
\item Source T1-A5. Question: Draw the structures of E, F, and G. Policy: derive\_in\_answer\_blind\_run\_or\_use\_problem\_stated\_fallback
\end{itemize}

\paragraph{Formalization target.}
create a compiling Lean file with sorry bodies at `IChO2026Problems/problem\_icho\_2026\_t1\_a6.lean`.
The Lean declarations must preserve every mathematical object, quantifier, hypothesis, side condition, approximation/error guarantee, and requested conclusion from the source statement.
Use Mathlib/Physlib/CRNT names found through LeanExplore where available. Prefer the configured benchmark Base library; introduce a local definition only when the source genuinely requires an object that the environment does not expose.
This is an answer-blind solve-phase task. Derive a candidate result only from the problem-side evidence above; do not assume a target value, select a tolerance around a desired value, or encode a desired result into a candidate domain.
For numerical questions, define the raw end-to-end quantity and apply an explicit source-derived rounding rule. For identification questions, characterize or prove uniqueness before recording the derived candidate.

\begin{theorem}[Icho Chemistry formalization target]
\label{thm:physics:icho_2026_t1_a6:target}
The assigned autoformalize agent should translate this IChO chemistry problem into Lean declarations in the covered file, with theorem and lemma proof bodies written as `by sorry`.
\end{theorem}
\begin{proof}
This is an autoformalization task, not a proof task. Produce faithful statements that can later be proved without weakening the source contract.
\end{proof}
% --- Archon physics formalization source end ---
